% Copyright 2026 Alan Szmyt. All rights reserved pending publication license.
\section{Introduction}
\label{sec:introduction}

Personalizing audio is not only a question of which music suits a person. It is
also a question of what that person wants to explore in a particular moment,
how the intended transition should unfold, what audio is actually produced and
heard, and how the person's response should---or should not---inform a later
proposal. Prior systems already connect parts of this problem. They select
music from current and desired affect while learning person-specific
associations \citep{janssen2012tune}; generate music along affective
trajectories inside physiological feedback loops
\citep{daly2016abcmi,ehrlich2019closedloop}; expose interpretable controls for
real-time affective generation \citep{agres2023affectmachine}; and mediate
estimated state through a semantic plan before controllable generation
\citep{zhang2026mindmelody}. These precedents make a broad claim to
``personalized,'' ``closed-loop,'' or ``interpretable'' music generation
untenable.

They also expose a more specific research problem. A requested transition is
not the same record as an estimated present state; an accepted semantic plan is
not the audio a model realizes; a generated artifact is not evidence that it
was played; and the affect a listener perceives in music is not necessarily the
affect that listener feels. Emotional intensity, immediate usefulness or harm,
and a later aftereffect are distinct again. If these transformations are
collapsed into a prompt, a model output, or a single success score, it becomes
difficult to reconstruct which consented context informed a proposal, where a
mismatch arose, what the person actually encountered, or what evidence could
justify changing a later proposal. We therefore treat inspectability of the
complete transformation chain as the research object; we do not treat
inspectability as a guarantee of correctness, meaningful consent, privacy,
safety, or benefit.

Antidote originated in one researcher's self-report of a strongly cathartic
intentional generative-audio session. That lived observation motivates asking
whether more deliberate person-and-moment audio journeys can be represented
and studied. It is not a study result, an efficacy signal, a neurological
mechanism, or evidence that the experience will generalize. Throughout this
manuscript it remains a hypothesis-generating observation.

We organize the resulting design space around four linked but
non-interchangeable terms. The \emph{person} contributes explicit authority,
semantic language, exclusions, and potentially relevant prior records. The
\emph{moment} is a consent-scoped and correctable representation of present
context, not an objective diagnosis or inferred need. The \emph{journey} is an
accepted, time-indexed progression toward an intended experience, not proof
that a state transition occurred. The \emph{audio} is the realized artifact,
whose identity and measured properties remain separate from its requested
meaning and from the person's response. Figure~\ref{fig:semantic-acoustic-response-loop}
introduces this person--moment--journey--audio chain. Its return path is a
hypothesis for later within-person inquiry, not an observed adaptive mechanism.

\AntidoteFigure{semantic-acoustic-response-loop}

\subsection{Research gap}
\label{sec:research-gap}

The relevant literature does not leave an empty technical category for
Antidote to fill. Longitudinal person-specific modeling, current-to-target
selection, generated affective trajectories, smooth real-time feedback,
interpretable musical controls, and semantic intervention planning all have
documented precedent \citep{janssen2012tune,daly2016abcmi,
ehrlich2019closedloop,agres2023affectmachine,zhang2026mindmelody}.
Furthermore, a recent preliminary brain--computer music study reported
nonsignificant target-emotion and time effects while person and trial variation
dominated its observed variance \citep{monroy2026minimalist}. Technical loop
closure therefore cannot be assumed to establish reliable affective change.

Instead, our bounded review identifies an operational gap: it did not find one
evaluated system that preserves, as a complete governed chain, (1) the context
a person permitted for a particular purpose, (2) user-extensible semantic
intent and exclusions, (3) an accepted time-indexed journey plan, (4) the
generated artifact and measured acoustic realization, (5) actual exposure,
(6) perceived expression, felt response, immediate usefulness or harm, and
later aftereffect as separate observations, and (7) any proposed within-person
update with its provenance and explicit human approval. ``Did not find'' is a
bounded-review result, not proof that no such system exists. The proposed
combination may constitute a useful research-system contribution; its global
novelty, technical feasibility, learnability, and benefit remain unresolved.

This framing shifts the central question away from whether a system can infer
an emotion and generate music. It asks whether a research instrument can make
the entire person-and-moment audio hypothesis inspectable enough to specify,
challenge, reproduce, and evaluate without promoting one record type into
another.

\subsection{Staged research questions}
\label{sec:staged-research-questions}

The manuscript separates a design question from two later feasibility
questions:

\begin{description}
  \item[RQ1---design.] How can a local research instrument represent,
    generate, and evaluate a person-and-moment audio journey while keeping
    consented context, semantic intent, planned controls, realized acoustics,
    exposure, response, and longitudinal updates inspectable and
    non-interchangeable?

  \item[RQ2---technical feasibility.] Can the implemented instrument preserve
    its declared authority, provenance, control-adherence, continuity,
    interruption, and recovery contracts across a complete local session?

  \item[RQ3---within-person feasibility.] Across repeated sessions, do
    provenance-linked person-and-moment records support useful advisory
    relationships among semantic intent, acoustic realization, and
    self-reported response without increasing burden, surprise, mismatch, or
    harm?
\end{description}

The design/protocol manuscript is structured to address RQ1 by formalizing the
research instrument and its evaluation boundary. RQ2 requires an end-to-end
implemented session and auditable technical checks. RQ3 additionally requires
a frozen protocol, resolved ethics and privacy gates, qualifying
repeated-session observations, and uncertainty-aware analysis. Neither later
question is answered by the originating experience, the literature review, a
repository fixture, or the proposed architecture.

\subsection{Bounded contributions}
\label{sec:introduction-contributions}

Within those limits, the complete design/protocol manuscript is scoped to five
contributions at explicitly different evidence states:

\begin{enumerate}
  \item it \emph{bounds} the problem against prior closed-loop,
    current-to-target, controllable, physiologically adaptive, and longitudinal
    music systems rather than presenting those ingredients as new;
  \item it \emph{formalizes} an end-to-end inspectability contract that keeps
    consented context, semantic intent, planning, acoustic realization,
    exposure, response layers, provenance, and advisory updating distinct;
  \item it \emph{proposes} a two-rate adaptive-audio architecture that
    separates uncertain, person-correctable deliberation from deterministic
    generate-ahead rendering, without claiming a validated controller or
    seamless listening experience;
  \item it \emph{requires} a local-first authority model in which context use,
    plan acceptance, generation, playback, retention, export, and
    personal-model updates remain interruptible and human-correctable, while
    distinguishing the partially implemented mock path from proposed runtime
    behavior; and
  \item it will \emph{specify prospectively} an evidence-disciplined
    feasibility boundary that retains null, negative, missing, interrupted, and
    adverse observations and never substitutes intended control, acoustic
    adherence, perceived expression, felt response, usefulness or harm, or
    aftereffect for one another.
\end{enumerate}

The remainder of the manuscript locates these contributions in prior work,
formalizes the system and its mathematical objects, specifies the prospective
feasibility method and results boundary, discusses only evidence-proportional
interpretations, states limitations and ethical constraints, and records
availability and reproducibility conditions before concluding with the gates
required to revisit RQ2 and RQ3.
